\chapter{Módulo 2 --- Estructuras de datos}

\begin{itemize}
\tightlist
\item
  \textbf{Duración estimada:} 1.5 h teoría + 2.5 h práctica
\item
  \textbf{Objetivos de aprendizaje.} Al finalizar, la persona podrá:

  \begin{enumerate}
  \def\labelenumi{\arabic{enumi}.}
  \tightlist
  \item
    \textbf{Crear} vectores, factores, listas, matrices y data frames.
  \item
    \textbf{Explicar} la coerción de tipos y el reciclaje de vectores.
  \item
    \textbf{Aplicar} indexación por posición, nombre y condición lógica
    (\passthrough{\lstinline![ ]!}, \passthrough{\lstinline![[ ]]!},
    \passthrough{\lstinline!$!}).
  \item
    \textbf{Clasificar} qué estructura conviene para un problema dado.
  \end{enumerate}
\end{itemize}

\section{Contenidos}

\begin{enumerate}
\def\labelenumi{\arabic{enumi}.}
\tightlist
\item
  Vectores atómicos: \passthrough{\lstinline!c()!},
  \passthrough{\lstinline!seq()!}, \passthrough{\lstinline!rep()!},
  \passthrough{\lstinline!length()!}.
\item
  Operaciones vectorizadas y reciclaje.
\item
  Coerción: jerarquía
  \passthrough{\lstinline!logical < integer < numeric < character!}.
\item
  Valores especiales: \passthrough{\lstinline!NA!},
  \passthrough{\lstinline!NULL!}, \passthrough{\lstinline!NaN!},
  \passthrough{\lstinline!Inf!}.
\item
  Factores: niveles y orden.
\item
  Matrices: \passthrough{\lstinline!matrix()!},
  \passthrough{\lstinline!dim()!}.
\item
  Listas: estructuras heterogéneas.
\item
  Data frames: la tabla de datos; \passthrough{\lstinline!str()!},
  \passthrough{\lstinline!head()!}, \passthrough{\lstinline!summary()!},
  \passthrough{\lstinline!nrow()!}.
\item
  Indexación: \passthrough{\lstinline![ ]!},
  \passthrough{\lstinline![[ ]]!}, \passthrough{\lstinline!$!} y filtros
  lógicos.
\end{enumerate}

\section{Desarrollo}

\subsection{Vectores y vectorización}

\begin{lstlisting}[language=R]
unidades <- c(12, 7, 20, 3, 15)
precios  <- c(4999, 7999, 899, 6499, 2499)
ingresos <- unidades * precios          # operación elemento a elemento
sum(ingresos)
mean(unidades)
length(unidades)
1:10 * 2                                # reciclaje de un escalar
\end{lstlisting}

\subsection{Coerción y NA}

\begin{lstlisting}[language=R]
c(1, "dos", TRUE)       # todo se vuelve character
c(1, TRUE, FALSE)       # 1 1 0 (logical -> numeric)
x <- c(4, NA, 10)
sum(x)                  # NA
sum(x, na.rm = TRUE)    # 14
is.na(x)                # FALSE TRUE FALSE
\end{lstlisting}

\subsection{Factores}

\begin{lstlisting}[language=R]
canal <- factor(c("En línea", "Departamental", "En línea", "Autoservicio"))
levels(canal)                   # orden alfabético por defecto
table(canal)
talla <- factor(c("M", "S", "L", "M"), levels = c("S", "M", "L"), ordered = TRUE)
talla < "L"
\end{lstlisting}

\subsection{Matrices y listas}

\begin{lstlisting}[language=R]
m <- matrix(1:6, nrow = 2)
dim(m)
m[2, 3]                         # fila 2, columna 3

tienda <- list(id = "T01", region = "Norte", ventas = c(120, 98, 143))
tienda$region
tienda[["ventas"]][2]           # 98
str(tienda)
\end{lstlisting}

\subsection{Data frames}

\begin{lstlisting}[language=R]
df <- data.frame(
  producto = c("Smartphone A", "Tablet", "Audífonos"),
  precio   = c(4999, 6499, 899),
  unidades = c(12, 4, 20)
)
str(df)
df$ingreso <- df$precio * df$unidades   # nueva columna
df[df$unidades > 10, ]                  # filas que cumplen condición
df[, c("producto", "ingreso")]          # columnas por nombre
nrow(df); ncol(df)
\end{lstlisting}

\subsection{Leer el conjunto de datos del curso (R base)}

\begin{lstlisting}[language=R]
ventas <- read.csv("datos/ventas.csv")
head(ventas)
summary(ventas$unidades)
\end{lstlisting}

\section{Actividades y ejercicios}

\begin{enumerate}
\def\labelenumi{\arabic{enumi}.}
\tightlist
\item
  Crea un vector con las temperaturas
  \passthrough{\lstinline!c(22, 25, NA, 30, 28)!} y calcula su promedio
  ignorando el \passthrough{\lstinline!NA!}.
\item
  Con \passthrough{\lstinline!ventas!}, ¿cuántas filas y columnas tiene?
  ¿Cuántos \passthrough{\lstinline!NA!} hay en
  \passthrough{\lstinline!unidades!}?
\item
  Convierte \passthrough{\lstinline!ventas$producto!} en factor y
  muestra cuántos registros hay por producto.
\item
  Extrae las filas de \passthrough{\lstinline!ventas!} del producto
  \passthrough{\lstinline!"Tablet"!} con más de 8 unidades.
\item
  Crea una lista \passthrough{\lstinline!resumen!} con el total de
  unidades y el producto más frecuente.
\end{enumerate}

\subsection{Soluciones}

\begin{lstlisting}[language=R]
# 1
temp <- c(22, 25, NA, 30, 28)
mean(temp, na.rm = TRUE)                       # 26.25

# 2
dim(ventas)
sum(is.na(ventas$unidades))                    # 40

# 3
ventas$producto <- factor(ventas$producto)
table(ventas$producto)

# 4
tablets <- ventas[ventas$producto == "Tablet" &
                  !is.na(ventas$unidades) &
                  ventas$unidades > 8, ]
head(tablets)

# 5
resumen <- list(
  total_unidades = sum(ventas$unidades, na.rm = TRUE),
  mas_frecuente  = names(which.max(table(ventas$producto)))
)
str(resumen)
\end{lstlisting}

\section{Evaluación (formativa)}

\begin{itemize}
\tightlist
\item
  \textbf{Quiz:} ¿Qué devuelve \passthrough{\lstinline!c(1, "a")!}?
  ¿Cuál es la diferencia entre \passthrough{\lstinline!lista["x"]!} y
  \passthrough{\lstinline!lista[["x"]]!}? ¿Para qué sirve
  \passthrough{\lstinline!levels=!} en
  \passthrough{\lstinline!factor()!}?
\item
  \textbf{Práctica calificada:} script que lea
  \passthrough{\lstinline!ventas.csv!}, cree la columna
  \passthrough{\lstinline!ingreso!} y muestre las 5 filas con mayor
  ingreso (pista: \passthrough{\lstinline!order()!}, con
  \passthrough{\lstinline!decreasing = TRUE!}).
\end{itemize}

\section{Referencias verificadas}

\begin{itemize}
\tightlist
\item
  \textbf{Libro:} Grolemund, G. (2014). \emph{Hands-On Programming with
  R}. O'Reilly Media. ISBN 978-1-4493-5901-0.
  \url{https://rstudio-education.github.io/hopr/} {[}EN{]} {[}OA en
  línea{]} --- cap. ``R Objects'' y ``R Notation''.
\item
  \textbf{Libro:} Wickham, H. (2019). \emph{Advanced R} (2.ª ed.). CRC
  Press. ISBN 978-0-8153-8457-1. \url{https://adv-r.hadley.nz/} {[}EN{]}
  {[}OA en línea, CC BY-NC-SA 4.0{]} --- cap. 3 ``Vectors'' y cap. 4
  ``Subsetting'' (lectura de profundización).
\item
  \textbf{Web:} R Core Team. \emph{An Introduction to R} --- secciones
  ``Simple manipulations; numbers and vectors'' y ``Lists and data
  frames''.
  \url{https://cran.r-project.org/doc/manuals/r-release/R-intro.html}
  {[}EN{]} {[}OA{]}
\item
  \textbf{Video:} freeCodeCamp.org. \emph{R Programming Tutorial --
  Learn the Basics of Statistical Computing} {[}Video{]}. YouTube.
  \url{https://youtu.be/_V8eKsto3Ug} {[}EN{]} --- secciones ``Data
  Formats'', ``Factors'', ``Entering Data'', ``Importing Data''.
\end{itemize}
